\section{Related Work}

\subsection{Financial peer groups and clustering}

Industry classifications provide a conventional way to identify comparable companies, but they do not fully describe financial similarity. Bhojraj et al. found differences in how well classification schemes group firms with similar returns, valuation multiples and financial ratios \citep{bhojraj2003}. Hoberg and Phillips used company product descriptions to develop industries that change over time \citep{hoberg2016}. Lee et al. identified peers through investors' co-search behaviour, revealing relationships across standard industry boundaries \citep{lee2015}. These studies support treating sectors as one description of similarity rather than the only one.

Clustering offers an alternative based on measured characteristics and a chosen distance measure \citep{greene2008}. Barradas et al. used PCA and K-Means to identify stock-market conditions associated with volatility \citep{barradas2023}. Chen et al. found that return-based commodity clusters aligned with market categories, whereas volatility-based clusters identified distress periods \citep{chen2021}. Le\'{o}n et al. used agglomerative clustering to examine changes in relationships between national equity-market indices following the global financial crisis \citep{leon2017}. Nourahmadi et al. compared K-Means, hierarchical clustering and affinity propagation for companies listed on the Tehran Stock Exchange \citep{nourahmadi2022}. These applications motivate comparing methods and features rather than assuming a single financial grouping.

Other studies use clusters for portfolio selection and trading. Lim and Ong argued that price-pattern groups can supplement industry-based diversification \citep{lim2021}. Cuomo et al. developed an unsupervised framework for market-neutral portfolios \citep{cuomo2022}, while Unlu and Xanthopoulos combined ensemble clustering with portfolio theory to reduce clustering variance \citep{unlu2021}. Han et al. applied clustering to pairs trading \citep{han2023}, and Nair et al. used stock-price clusters for trading recommendations \citep{nair2017}. These applications illustrate possible financial uses. The present study evaluates the groups themselves, including whether their members have similar subsequent return movements, rather than testing a trading strategy.

\subsection{Representation learning and cluster validation}

Dimensionality reduction represents observations with fewer coordinates, helping readers inspect patterns across several variables. UMAP constructs a lower-dimensional representation using relationships between nearby observations \citep{mcinnes2018}. t-SNE offers another way to visualise relationships in high-dimensional data \citep{vandermaaten2008}. However, clear separation in a two-dimensional map does not necessarily imply equally strong separation in the original features. The map and the financial variables therefore provide different views of the same assignments.

Internal validation measures assess the closeness of observations within groups and separation between groups. The Silhouette Score compares an observation's distance from its own group with its distance from the nearest alternative group \citep{rousseeuw1987}. The Davies-Bouldin Index compares within-group spread with separation between groups \citep{davies1979}. Higher silhouette and lower Davies-Bouldin values indicate better separation under their respective definitions. Comparative studies show that different validation measures can favour different groupings \citep{arbelaitz2013,vendramin2010}.

Evaluating only the reduced map can therefore give an incomplete picture. We assess the same assignments in both the map and the original scaled financial features. Differences between these scores describe separation in each space; they do not directly measure geometric distortion or establish that neighbourhoods were preserved. Comparisons across mapping methods, features, cluster counts and random seeds assess how much the findings depend on modelling choices.

\subsection{Tracking clusters over time}

Repeated clustering raises questions that do not arise when analysing a single dataset. Groups may gain or lose members, split, merge, disappear or emerge. Spiliopoulou et al. developed MONIC to monitor these types of cluster transitions \citep{spiliopoulou2006}. Greene et al. used Jaccard similarity to match communities with their predecessors in dynamic networks \citep{greene2010}.

An alternative is to build temporal consistency into the clustering objective. Chakrabarti et al. developed evolutionary clustering to balance the fit to current data with agreement with earlier groupings \citep{chakrabarti2006}. This differs from fitting each date independently and then comparing the results. Independent fitting allows changes in observed peer relationships to remain visible rather than explicitly favouring continuity.

The present study uses independent fitting followed by membership matching. Jaccard similarity measures overlap between groups, and Hungarian assignment finds the one-to-one matching with the highest total similarity \citep{greene2010,kuhn1955}. Matching makes it possible to track groups despite arbitrary cluster numbers. One-to-one matches can nevertheless simplify splits and mergers. We therefore examine overlap matrices and individual histories alongside counts of label changes, and distinguish index entries and exits from changes among continuing securities.

\subsection{Interpretability and anomaly screening}

Unusual financial profiles can reflect liquidity constraints, changing operations, capital structure or other company characteristics. They are not automatically errors or unfavourable investments. Earlier studies have examined anomaly detection in financial applications, including fraud detection \citep{ahmed2016} and unusual financial-market behaviour \citep{ahmed2017,cheong2021}. Their definitions of anomalies differ from the financial-profile screening used here. These studies motivate investigating unusual observations; they do not establish that our cross-sectional Mahalanobis flags identify economically significant anomalies.

Clusters formed in a nonlinear map are not expressed directly as rules on the original features. SHAP assigns feature contributions to outputs from a fitted model \citep{lundberg2017,merrick2020,rozemberczki2022}. A surrogate classifier can learn to reproduce cluster labels from the financial variables. SHAP then describes how that classifier uses the variables to reproduce the labels.

These explanations concern the surrogate, not the clustering algorithm itself. Their usefulness depends on surrogate fidelity: how closely the classifier reproduces the assignments, including for instruments held out during training. Direct feature profiles offer a separate basis for interpretation when fidelity is weak or cannot be estimated. Neither surrogate attributions nor measured feature changes establish a causal explanation for a company's movement between groups.

\subsection{External validation and research gap}

Internal scores describe groups in the features and representation used to construct them. They do not establish whether those groups capture relationships in observations excluded from fitting. Financial applications of clustering motivate testing this distinction rather than inferring practical usefulness from visual separation alone \citep{lim2021,cuomo2022,han2023}.

Here, subsequent return co-movement provides an external comparison for the June and December 2025 groupings. We test whether cluster peers have more strongly correlated daily returns over the following six months than sector peers. Sectors are a practical comparator, not ground-truth cluster labels. This comparison addresses return similarity; it does not establish trading profitability or validate individual anomaly flags.

The evaluation problem is therefore broader than choosing the method with the highest internal score. Separation in the map and financial features, agreement across model settings, membership changes, surrogate explanations and subsequent returns answer different questions. The reviewed studies motivate these perspectives, but do not make them interchangeable. Our empirical contribution is to examine them together in repeated JSE Top 40 cross-sections, including cases where they support different interpretations. The workflow combines established methods; it introduces neither a new clustering algorithm nor a return-prediction strategy.
